\documentclass[12 pt, letterpaper]{hmcpset}
\usepackage{hw-preamble}

\name{Jayden Thadani}
\class{MATH 171 --- Abstract Algebra I}
\assignment{Homework 13}
\duedate{5th April 2024}

\begin{document}

\begin{problem}[Dummit \& Foote \S7.1 \#14]
	Let $x$  be a nilpotent element of the commutative ring $R$ 
	\begin{enumerate}
		\item Prove that $x$ is either zero or a zero divisor.
		\item Prove that $rx$ is nilpotent for all $r \in R$.
		\item Prove that $1 + x$ is a unit in $R$.
		\item Deduce that the sum of a nilpotent element and a unit is a unit.
	\end{enumerate}
\end{problem}

\begin{solution}
	\begin{enumerate}
		\item Suppose $x$ is nonzero. Then there exists some least $n \in \mathbb{N}$ such that $x^{n} = 0$ with $n > 1$. Thus, $x x^{n - 1} = 0$ and $x^{n - 1}$ is nonzero (since $n$ is the least positive integer that gives $x^{n} = 0$). That is, $x$ is a zero divisor.
		\item Suppose $x^{n} = 0$. Then, for any $r \in R$, we have
			\[
				\begin{split}
					(rx)^{n} & = \underbrace{(rx) \dots (rx)}_{n \text{ times}} \\
						 & = r^{n} x^{n} \\
						 & = r^{n} 0 \\
						 & = 0.
				\end{split}
			\]
			where the second step follows from the commutativity of $R$. Thus, $rx$ is nilpotent.
		\item Suppose $x^{n} = 0$. Note that the polynomial $1 - y^{n} \in \mathbb{Z}[x]$ has root $y = - 1$ if $n$ is even and the polynomial $1 + y^{n}$ has root $y = -1$ if $n$ is odd. Specifically, if $n$ is even, we can write
			\[
				\begin{split}
					1 - y^{n} & = (1 - y)(y^{n - 1} + y^{n - 2} + \dots + y + 1) \\
						  & = (1 - y)(1 + y)(1 + y^{2} + \dots + y^{n - 2})
				\end{split}
			\]
			and
			\[
				1 + y^{n} = (1 + y)(y^{n - 1} - y^{n - 2} + \dots + y^{2} - y + 1).
			\] \\
			Since multiplication distributes over addition just as well in $R$ as it does in $\mathbb{Z}[x]$, we can write
			\[
				\begin{split}
					(1 + x)(1 - x)(1 + x^{2} + \dots + x^{n - 2}) & = 1 - x^{n} \\
										      & = 1 - 0 \\
										      & = 1
				\end{split}
			\]
			for $n$ even and
			\[
				\begin{split}
					(1 + x)(x^{n - 1} - x^{n - 2} + \dots + x^{2} - x + 1) & = 1 + x^{n} \\
											       & = 1 + 0 \\
											       & = 1
				\end{split}
			\]
			for $n$ odd. Thus, we have that $1 + x$ is a unit in both cases.
		\item We can use a similar proof to show that for any unit $u \in R$, $u + x$ is a unit. \\\\
			First note that for any $n \in \mathbb{N}$, $u^{n}$ is a unit --- just consider $u^{n} (u^{-1})^{n} = 1$. Also note that for any unit $u$, if $a, b \in R$ give $ab = u$, then $a, b$ must be units too. We can see this by first noticing if one is a unit, then the other must be. Without loss of generality, suppose $a$ is a unit. Then $b = a^{-1} u$ must be a unit since $u^{-1} a a^{-1} u = 1$. Notice that $a$ must be a unit since $a (b u^{-1}) = 1$, and thus both must be units. \\\\
			Then by the same factorisation trick, for $n$ such that $x^{n} = 0$, we can write $u^{n} + x^{n}$ or $u^{n} - x^{n}$ (depending on whether $n$ is odd or even) as $(u + x) p(x)$ for some $p \in \mathbb{Z}[x]$. We can use the canonical homomorphism of $\mathbb{Z}$ into every ring to replace the coefficients $a \in \mathbb{Z}$. That is, use $\varphi : \mathbb{Z} \to R$ by $1 \mapsto 1_{R}$ to write $a \mapsto \underbrace{1 + \dots + 1}_{a \text{ times}}$. \\
			Since distributivity still works, we have for all $n \in \mathbb{N}$
			\[
				(u + x) r = u^{n}
			\]
			for some $r \in R$. Then, by our previous results that $u^{n}$ is a unit and that two elements that multiply to form a unit must be a unit, we have that $u + x$ is a unit.
	\end{enumerate}
\end{solution}

\pagebreak

\begin{problem}[Dummit \& Foote \S7.2 \#10]
	Consider the following elements of the integral group ring $\mathbb{Z} S_{3}$ :
	\[
		\alpha = 3 (1 \, 2) - 5 (2 \, 3) + 14 (1 \, 2 \, 3) \quad \beta = 6 (1) + 2 (2 \, 3) - 7 (1 \, 3 \, 2)
	\]
	(where $(1)$ is the identity of $S_{3}$) . Compute the following elements:
	\begin{enumerate}
		\item $\alpha + \beta$
		\item $2 \alpha - 3 \beta$
		\item $\alpha \beta$
		\item $\beta \alpha$
		\item $\alpha^{2}$
	\end{enumerate}
\end{problem}

\pagebreak

\begin{problem}[Dummit \& Foote \S7.4 \#7]
	Let $R$ be a commutative ring with $1$. Prove that the principal ideal generated by $x$ in the polynomial ring $R[x]$ is a prime ideal if and only if $R$ is an integral domain. Prove that $(x)$ is a maximal ideal if and only if $R$ is a field.
\end{problem}

\pagebreak

\begin{problem}[Dummit \& Foote \S7.4 \#8]
	Let $R$ be an integral domain. Prove that $(a) = (b)$ for some elements $a, b \in R$ if and only if $a = ub$ for some unit $u$ of $R$.
\end{problem}

\pagebreak

\begin{problem}[Dummit \& Foote \S7.4 \#11]
	Assume that $R$ is commutative. Let $I$ and $J$ be ideals of $R$ and assume $P$ is a prime ideal of $R$ that contains $IJ$ (for example, if $P$ contains $I \cap J$). Prove either $I$ or $J$ is containe din $P$.
\end{problem}

\end{document}
